\chapter{Mayer 团簇展开与位力方程}
\label{chap:mayer-virial}

巨配分函数 \(\Xi\) 把所有粒子数的正则配分函数汇在一起，而它的对数 \(\log\Xi\) 会精确消去彼此不连通的贡献。这个简单事实是低密度展开的核心：经典气体中它表现为连通 Mayer 图，量子气体中则表现为正则配分函数的累积量。把所得级数从逸度反演到粒子数密度，便得到位力状态方程。热波长、质心质量和内部简并度必须从一开始固定，否则第二位力系数很容易因归一化约定不同而相差一个整体因子。

\section{Mayer 展开与量子团簇的归一化}

对质量为 \(m\) 的三维单组分气体定义热德布罗意波长
\begin{equation}
 \lambda_T(m)
 =\sqrt{\frac{2\pi\hbar^2\beta}{m}},
 \qquad
 \lambda_T:=\lambda_T(m).
 \label{eq:phys-sm-thermal-wavelength}
\end{equation}
单组分逸度仍为 \(z=\ee^{\beta\mu}\)。无量纲量子团簇系数 \(b_n\) 定义为
\begin{equation}
 \frac{1}{\mathcal V}\log\Xi(z)
 =\frac{1}{\lambda_T^3}
 \sum_{n=1}^{\infty}b_n z^n.
 \label{eq:phys-sm-cluster-convention}
\end{equation}
因此 \(b_n\) 无量纲。后面出现的位力系数 \(B_n\) 则以粒子数密度为变量，量纲为 \(\mathrm{length}^{3(n-1)}\)；两套系数不能混称为同一个“第二位力系数”。

先从经典可积稳定势说明连通图为何进入 \(\log\Xi\)。这一结构随后迁移到量子正则累积量；二者共享的是“对数抽取连通贡献”，而不是相同的微观积分。

经典对势气体提供最透明的连通组合结构原型。设
\[
 U_N(\mathbf r_1,\ldots,\mathbf r_N)
 =\sum_{1\le i<j\le N}u(\mathbf r_i-\mathbf r_j).
\]
假设势稳定，即存在 \(B_{\rm stab}\ge0\) 使任意 \(N\) 和任意构型都满足
\begin{equation}
 U_N\ge-B_{\rm stab}N,
 \label{eq:phys-sm-stable-potential}
\end{equation}
并假设 Mayer 键可积。定义
\begin{equation}
 \mathfrak f_{ij}=\ee^{-\beta u(\mathbf r_i-\mathbf r_j)}-1,
 \qquad
 \mathcal C_{\rm M}(\beta)=\int_{\mathbb R^3}|\mathfrak f(\mathbf r)|\dd^3r<\infty.
 \label{eq:phys-sm-mayer-f}
\end{equation}
用有量纲经典活度 \(\zeta\) 写巨配分函数
\begin{equation}
 \Xi_{\rm cl}(\zeta)
 =\sum_{N=0}^{\infty}\frac{\zeta^N}{N!}
 \int_{\mathcal V^N}
 \prod_{i=1}^N\dd^3r_i
 \prod_{i<j}(1+\mathfrak f_{ij}).
 \label{eq:phys-sm-mayer-grand}
\end{equation}
展开最后的乘积，任意标号图 \(G\) on \(\{1,\ldots,N\}\) 贡献所有所选边的 \(\mathfrak f_{ij}\) 乘积。指数公式说明取对数以后只留下连通图：
\begin{equation}
 \log\Xi_{\rm cl}
 =\sum_{n=1}^{\infty}\frac{\zeta^n}{n!}
 \int_{\mathcal V^n}\prod_{i=1}^n\dd^3r_i
 \sum_{G\in\mathcal C_n}
 \prod_{(ij)\in\mathsf E(G)}\mathfrak f_{ij},
 \label{eq:phys-sm-general-mayer}
\end{equation}
其中 \(\mathcal C_n\) 是 \(n\) 个标号顶点上的连通简单图。

\begin{proof}[解]
要说明对配分函数 \(\Xi\) 取对数为何只留下连通图，先把图按连通块分解，再写连通权重的指数生成函数，最后取对数。

图的连通块分解。任意标号图 \(G\) on \(\{1,\ldots,N\}\) 唯一分解为若干连通块 \(G_1\sqcup\cdots\sqcup G_k\)，各块顶点集互斥且无跨块边。若块 \(G_i\) 有 \(n_i\) 个顶点，则 \(\sum_i n_i=N\)，且 Mayer 键乘积因子化：
\[
 \prod_{(j,l)\in\mathsf E(G)}\mathfrak f_{jl}
 =\prod_{i=1}^k\prod_{(j,l)\in\mathsf E(G_i)}\mathfrak f_{jl}.
\]
对位置的积分也因此因子化为各块独立积分。这里 \(\mathfrak f_{jl}=\ee^{-\beta V(\mathbf r_j-\mathbf r_l)}-1\) 是 Mayer \(f\)-函数，满足 \(\mathfrak f_{jl}\to0\)（\(|\mathbf r_j-\mathbf r_l|\to\infty\)），保证短程势下积分收敛。

连通块权重的指数生成函数。令 \(W_n\) 表示所有 \(n\)-点连通图的积分权重和：
\[
 W_n=\int_{\mathcal V^n}\prod_{i=1}^n\dd^3r_i
 \sum_{G\in\mathcal C_n}
 \prod_{(j,l)\in\mathsf E(G)}\mathfrak f_{jl}.
\]
一个一般 \(N\)-点图由连通块 \(G_1,\ldots,G_k\) 组成，块大小 \(n_1,\ldots,n_k\) 满足 \(\sum n_i=N\)。\(N\) 个 label 分配到各块有 \(N!/(n_1!\cdots n_k!)\) 种方式，但块内部已按标号图求和，且块之间无顺序。因此全部图（连通与非连通）的配分函数为
\[
 \Xi_{\rm cl}
 =\sum_{N=0}^{\infty}\frac{\zeta^N}{N!}
 \sum_{\substack{k\ge0\\n_1+\cdots+n_k=N}}
 \frac{N!}{n_1!\cdots n_k!}
 \prod_{i=1}^k W_{n_i}
 =\prod_{n=1}^{\infty}\sum_{k_n=0}^{\infty}
 \frac{1}{k_n!}\left(\frac{\zeta^nW_n}{n!}\right)^{k_n}
 =\exp\!\left(\sum_{n=1}^{\infty}\frac{\zeta^n}{n!}W_n\right),
\]
其中 \(k_n\) 是大小为 \(n\) 的连通块个数，最后一步用了 \(\sum_{k=0}^\infty x^k/k!=\ee^x\)。这一步的代数本质是指数公式（指数公式）：集合分拆的指数生成函数等于其连通分量的指数生成函数的指数。

取对数。对上式取对数立即得到\eqnrefs{eq:phys-sm-general-mayer}：
\[
 \log\Xi_{\rm cl}
 =\sum_{n=1}^{\infty}\frac{\zeta^n}{n!}W_n
 =\sum_{n=1}^{\infty}\frac{\zeta^n}{n!}
 \int_{\mathcal V^n}\prod_{i=1}^n\dd^3r_i
 \sum_{G\in\mathcal C_n}
 \prod_{(j,l)\in\mathsf E(G)}\mathfrak f_{jl}.
\]
非连通图并未被"忽略"：它们由连通块权重的指数自动重建。取对数的效果恰好剥离了这一指数结构，只留下连通扇区。

对 \(n=2\)，连通简单图只有一条边，\(\mathcal C_2=\{(12)\}\)，故
\begin{equation*}
 W_2=\int_{\mathcal V^2}\dd^3r_1\dd^3r_2\,\mathfrak f_{12}
 \xrightarrow{\mathcal V\to\infty}
 \mathcal V\int_{\mathbb R^3}\dd^3r\,\bigl(\ee^{-\beta V(r)}-1\bigr),
\end{equation*}
代入 \(\log\Xi=\zeta W_1+\zeta^2W_2/2+\cdots\) 并与位力展开比较，得经典第二位力系数
\begin{equation*}
 B_2^{\rm cl}=-\frac12\int_{\mathbb R^3}\dd^3r\,\bigl(\ee^{-\beta V(r)}-1\bigr),
\end{equation*}
即 Mayer 第二位力积分。对硬球势 \(V(r)=\infty\,(r<d)\)、\(0\,(r>d)\)，积分给出 \(B_2=2\pi d^3/3\)，即四倍硬球体积。

量子气体中 Mayer 键被正则迹 \(Q_n\) 替代，连通性由累积量结构保证：\(\log\Xi=\sum b_nz^n\) 中 \(b_n\) 是 \(Q_1,\ldots,Q_n\) 的多项式，恰好删除所有非连通分拆。此结构在量子场论中对应 Feynman 图的连通图定理：\(\log Z[J]\) 的生成泛函只含连通图，是统计力学与 QFT 共用的组合骨架。
\end{proof}

稳定性和可积性还给出团簇展开的一个充分收敛域。例如 Penrose--Ruelle 型估计可写成
\begin{equation}
 |\zeta|\mathcal C_{\rm M}(\beta)\ee^{2\beta B_{\rm stab}+1}<1.
 \label{eq:phys-sm-mayer-sufficient-condition}
\end{equation}
这个条件保证低活度解析性，但只是充分条件，不等于真实相变边界。

量子气体中不再使用经典 Mayer 键，而从已经施加交换统计的正则迹出发。巨配分函数的对数仍按累积量组合抽取连通部分，因此前三阶系数可以直接由 \(Q_1,Q_2,Q_3\) 得到。

量子体系固定粒子数的正则配分函数定义为
\begin{equation}
 Q_N=\Tr_{\mathscr H_N^{(\eta)}}\ee^{-\beta H_N},
 \qquad
 \eta=+1\ \text{(boson)},\quad
 \eta=-1\ \text{(fermion)},
 \label{eq:phys-sm-canonical-qn}
\end{equation}
其中 \(\mathscr H_N^{(\eta)}\) 已经是完全对称或反对称的 \(N\)-粒子子空间，因此不能再额外除以 \(N!\)。巨配分函数为
\begin{equation}
 \Xi(z)=\sum_{N=0}^{\infty}z^NQ_N,
 \qquad Q_0=1.
 \label{eq:phys-sm-quantum-grand}
\end{equation}
由 \(z\Xi'(z)=\Xi(z)\,z\partial_z\log\Xi(z)\) 可先得到任意阶正则配分函数与团簇系数之间的递推。把\eqnrefs{eq:phys-sm-quantum-grand} 与\eqnrefs{eq:phys-sm-cluster-convention} 代入并比较 \(z^N\) 系数，
\begin{equation}
 NQ_N=\frac{\mathcal V}{\lambda_T^3}
 \sum_{n=1}^{N}n b_n Q_{N-n},
 \qquad Q_0=1.
 \label{eq:phys-sm-general-cluster-recurrence}
\end{equation}
于是
\begin{equation}
 b_N=\frac{\lambda_T^3}{\mathcal V}Q_N
 -\frac1N\sum_{n=1}^{N-1}n b_n Q_{N-n}.
 \label{eq:phys-sm-cluster-recurrence-solved}
\end{equation}
这说明每个 \(b_N\) 都是正则迹中恰好去除所有低阶非连通分拆后的累积量。前三阶特殊化为
\begin{align}
 b_1&=\frac{\lambda_T^3}{\mathcal V}Q_1,
 \label{eq:phys-sm-b1-canonical}\\
 b_2&=\frac{\lambda_T^3}{\mathcal V}
 \left(Q_2-\frac{Q_1^2}{2}\right),
 \label{eq:phys-sm-b2-canonical}\\
 b_3&=\frac{\lambda_T^3}{\mathcal V}
 \left(Q_3-Q_1Q_2+\frac{Q_1^3}{3}\right).
 \label{eq:phys-sm-b3-canonical}
\end{align}
这些是巨-正则累积量。\(b_3\) 已经含真正的三体连通扇区，因此不能由二体 Beth--Uhlenbeck 相移公式简单延拓得到。

\begin{proof}[解]
巨配分函数的 Taylor 展开为
\begin{equation*}
 u(z)=zQ_1+z^2Q_2+z^3Q_3+O(z^4),
 \qquad \Xi=1+u.
\end{equation*}
累积量生成函数是 $\log\Xi=\log(1+u)$。对 $|u|<1$ 使用 Taylor 展开
\begin{equation*}
 \log(1+u)=u-\frac{u^2}{2}+\frac{u^3}{3}+O(u^4).
\end{equation*}
此展开的收敛半径为 1，对应 $|z|<z_0$ 的小逸度区；在相变处收敛半径触及奇点，累积量展开失效。

由 $u(z)$ 的展开式，
\begin{align*}
 u^2
 &=\bigl(zQ_1+z^2Q_2+O(z^3)\bigr)^2
 =z^2Q_1^2+2z^3Q_1Q_2+O(z^4),\\
 u^3
 &=\bigl(zQ_1+O(z^2)\bigr)^3
 =z^3Q_1^3+O(z^4).
\end{align*}
交叉项 $2z^3Q_1Q_2$ 来自 $(zQ_1)(z^2Q_2)+(z^2Q_2)(zQ_1)$。$u^3$ 中 $z^3Q_1^3$ 来自 $(zQ_1)^3$，其余项至少含一个 $O(z^2)$ 因子故为 $O(z^4)$。

把 $u=zQ_1+z^2Q_2+z^3Q_3$ 代入 $\log\Xi=\log(1+u)$，展开到 $z^3$ 阶：
\begin{align*}
 \log\Xi
 &=\bigl(zQ_1+z^2Q_2+z^3Q_3\bigr)
  -\frac{1}{2}\bigl(z^2Q_1^2+2z^3Q_1Q_2\bigr)
  +\frac{1}{3}\bigl(z^3Q_1^3\bigr)+O(z^4)\\
 &=zQ_1
  +z^2\left(Q_2-\frac{Q_1^2}{2}\right)
  +z^3\left(Q_3-Q_1Q_2+\frac{Q_1^3}{3}\right)
  +O(z^4).
\end{align*}

由\eqnrefs{eq:phys-sm-cluster-convention}，$\log\Xi=\sum_{n\ge1}b_nz^n$，逐次比较 $z$、$z^2$、$z^3$ 的系数即得到\eqnrefs{eq:phys-sm-b1-canonical}--\eqnrefs{eq:phys-sm-b3-canonical}。
\end{proof}

设单粒子内部 Hilbert 空间维数为 \(d_{\rm int}\)，且所有内部态等能。理想气体满足
\begin{equation}
 b_n^{(0)}
 =\frac{d_{\rm int}\,\eta^{n+1}}{n^{5/2}}.
 \label{eq:phys-sm-ideal-cluster}
\end{equation}
这里不从压力级数整体提出 \(d_{\rm int}\)。因此若另一约定定义“每个自旋分量的 \(\bar b_n\)”，两套系数会相差相应的整体内部简并度。

\begin{proof}[解]
理想巨配分函数。单粒子色散 $\varepsilon_{\mathbf p}=p^2/(2m)$，内部 Hilbert 空间维数 $d_{\rm int}$。理想巨配分函数（$\eta=+1$ Bose，$\eta=-1$ Fermi）满足
\begin{equation*}
 \log\Xi_0
 =-\eta d_{\rm int}
 \sum_{\mathbf p}
 \log\left(1-\eta z\,\ee^{-\beta\varepsilon_{\mathbf p}}\right).
\end{equation*}
此式来自巨配分函数的乘积形式 \(\Xi_0=\prod_{\mathbf p,\sigma}(1-\eta z\ee^{-\beta\varepsilon_{\mathbf p}})^{-\eta}\)，取对数后对内部态 \(\sigma\) 求和给出 \(d_{\rm int}\) 因子。Bose 情形 \(\eta=+1\) 要求 \(|z|<1\) 以保证收敛；Fermi 情形 \(\eta=-1\) 对任意 \(z\) 解析。

展开对数。利用 Taylor 展开
\begin{equation*}
 -\eta\log(1-\eta x)
 =\sum_{n=1}^{\infty}\frac{\eta^{n+1}}{n}\,x^n,
\end{equation*}
（Bose 时 $\eta=+1$，得 $-\log(1-x)=\sum x^n/n$；Fermi 时 $\eta=-1$，得 $\log(1+x)=\sum(-1)^{n+1}x^n/n$），代入第一步，
\begin{equation*}
 \log\Xi_0
 =d_{\rm int}\sum_{n=1}^{\infty}\frac{\eta^{n+1}}{n}\,z^n
 \sum_{\mathbf p}\ee^{-n\beta p^2/(2m)}.
\end{equation*}
注意指数中 \(n\beta\) 而非 \(\beta\)：第 \(n\) 项对应 \(n\) 个粒子在同一动量态的"环图"贡献，等效温度 \(T/n\)。

热力学极限。动量和在热力学极限下化为积分，
\begin{equation*}
 \sum_{\mathbf p}\ee^{-n\beta p^2/(2m)}
 \xrightarrow{\mathcal V\to\infty}
 \frac{\mathcal V}{(2\pi\hbar)^3}\int\ee^{-n\beta p^2/(2m)}\dd^3p
 =\frac{\mathcal V}{\lambda_T^3\,n^{3/2}},
\end{equation*}
其中 $\lambda_T=(2\pi\hbar^2\beta/m)^{1/2}$ 是热波长，$n^{-3/2}$ 来自三维 Gauss 积分
\begin{equation*}
 \int_{\mathbb R^3}\ee^{-n\beta p^2/(2m)}\dd^3p
 =\left(\frac{2\pi m}{n\beta}\right)^{3/2}
 =\frac{(2\pi\hbar)^3}{\lambda_T^3\,n^{3/2}}.
\end{equation*}

合并。
\begin{equation*}
 \log\Xi_0
 =\frac{\mathcal V}{\lambda_T^3}
 \sum_{n=1}^{\infty}
 \frac{d_{\rm int}\,\eta^{n+1}}{n^{5/2}}\,z^n.
\end{equation*}
与团簇展开 $\log\Xi=\sum_n b_n z^n$ 比较，$b_n=d_{\rm int}\,\eta^{n+1}/n^{5/2}$，即\eqnrefs{eq:phys-sm-ideal-cluster}。$n^{-5/2}=n^{-1}\cdot n^{-3/2}$ 中，$n^{-1}$ 来自对数展开，$n^{-3/2}$ 来自动量积分。

Bose 气体 \(\eta^{n+1}=+1\)，所有 \(b_n>0\)，对应有效吸引（聚束倾向）；Fermi 气体 \(\eta^{n+1}=(-1)^{n+1}\)，\(b_n\) 交替变号，对应有效排斥（Pauli 不相容）。第二位力系数 \(B_2=-\lambda_T^3 b_2/b_1^2=-\eta\,\lambda_T^3/(2^{5/2})\)：Bose 为负（吸引修正），Fermi 为正（排斥修正），与经典 \(B_2=0\) 的差异即交换效应。

若单粒子色散改为 \(\varepsilon_p=\sqrt{p^2c^2+m^2c^4}-mc^2\)，Gauss 积分被 Bessel 函数 \(K_2(n\beta mc^2)\) 替代，团簇系数不再是简单的 \(n^{-5/2}\) 而含相对论修正因子。在非相对论极限 \(\beta mc^2\gg1\) 下 \(K_2(x)\sim\sqrt{\pi/(2x)}\ee^{-x}\) 退化为经典结果。
\end{proof}

\section{位力展开与一般二体团簇}

粒子数密度记为 \(\varrho\)。由巨势
\begin{equation}
 \varrho
 =\frac{z}{\mathcal V}\frac{\partial\log\Xi}{\partial z}
 =\frac{1}{\lambda_T^3}
 \sum_{j\ge1}j b_jz^j.
 \label{eq:phys-sm-density-cluster}
\end{equation}
压力满足 \(\beta P=\mathcal V^{-1}\log\Xi\)。令 \(x=\varrho\lambda_T^3\)，则由\eqnrefs{eq:phys-sm-density-cluster} 有 \(x(z)=b_1z+O(z^2)\)。只要 \(b_1\neq0\)，解析逆函数定理保证原点附近存在唯一解析反函数 \(z=z(x)\)，因此位力级数并不是额外假设，而是团簇级数在局部可逆条件下的重参数化。位力展开以物理密度为变量：
\begin{equation}
 \beta P
 =\varrho+B_2\varrho^2+B_3\varrho^3+\cdots.
 \label{eq:phys-sm-virial-definition}
\end{equation}
逐阶反演逸度级数得到
\begin{align}
 B_2&=-\lambda_T^3\frac{b_2}{b_1^2},
 \label{eq:phys-sm-virial-b2}\\
 B_3&=\lambda_T^6
 \left(
 \frac{4b_2^2}{b_1^4}
 -\frac{2b_3}{b_1^3}
 \right).
 \label{eq:phys-sm-virial-b3}
\end{align}
因此 \(b_n\) 无量纲，而 \([B_n]=\text{length}^{3(n-1)}\)。

\begin{proof}[解]
定义 \(x=\varrho\lambda_T^3\)，则
\[
 x=b_1z+2b_2z^2+3b_3z^3+O(z^4).
\]
此式来自 \(\varrho=z\partial_z\log\Xi/\mathcal V=\lambda_T^{-3}\sum_{j\ge1}jb_jz^j\)，乘以 \(\lambda_T^3\) 即得 \(x\)。设
\[
 z=A_1x+A_2x^2+A_3x^3+O(x^4).
\]
逐阶比较给出
\[
 A_1=\frac1{b_1},
 \qquad
 A_2=-\frac{2b_2}{b_1^3},
 \qquad
 A_3=\frac{8b_2^2}{b_1^5}-\frac{3b_3}{b_1^4}.
\]
具体推导：把 \(z=A_1x+A_2x^2+A_3x^3\) 代入 \(x=b_1z+2b_2z^2+3b_3z^3\)，
\begin{align*}
 x &= b_1(A_1x+A_2x^2+A_3x^3)
 +2b_2(A_1^2x^2+2A_1A_2x^3)\\
 &\quad+3b_3A_1^3x^3+O(x^4),
\end{align*}
比较 \(x\) 系数：\(1=b_1A_1\)，得 \(A_1=1/b_1\)；比较 \(x^2\)：\(0=b_1A_2+2b_2A_1^2\)，得 \(A_2=-2b_2/b_1^3\)；比较 \(x^3\)：\(0=b_1A_3+4b_2A_1A_2+3b_3A_1^3\)，代入前两式得 \(A_3=8b_2^2/b_1^5-3b_3/b_1^4\)。代回
\[
 \beta P\lambda_T^3=b_1z+b_2z^2+b_3z^3+O(z^4)
\]
后，
\[
 \beta P\lambda_T^3
 =x-\frac{b_2}{b_1^2}x^2
 +\left(
 \frac{4b_2^2}{b_1^4}
 -\frac{2b_3}{b_1^3}
 \right)x^3+O(x^4),
\]
恢复 \(x=\varrho\lambda_T^3\) 即得\eqnrefs{eq:phys-sm-virial-b2}--\eqnrefs{eq:phys-sm-virial-b3}。

反演的解析性由隐函数定理保证：只要 \(b_1\neq0\)（理想气体 \(b_1=d_{\rm int}\neq0\) 满足），\(\varrho(z)\) 在 \(z=0\) 附近局部可逆，位力级数收敛半径由最近的团簇级数奇点决定。相变处位力级数发散，对应 \(\varrho(z)\) 的多值性或支点。
\end{proof}

位力系数只依赖低阶团簇数据。为了把二阶相互作用部分写成散射信息而不在人为的参考热波长上做选择，先使用有量纲的多组分连通系数，并完成质心--相对运动分离。若组分 \(a\) 的粒子数 \(N_a\) 各自守恒，才可以给每个守恒荷引入独立逸度 \(z_a\)，并写
\begin{equation}
 \frac1{\mathcal V}\log\Xi(\{z_a\})
 =\sum_{\mathbf n\neq0}
 \mathfrak c_{\mathbf n}
 \prod_a z_a^{n_a},
 \qquad
 [\mathfrak c_{\mathbf n}]=\text{length}^{-3}.
 \label{eq:phys-sm-multicomponent-cluster}
\end{equation}
若存在 Rabi 耦合、自旋转换或其他使单独 \(N_a\) 不守恒的项，就不能先赋予每个内部标签独立化学势；此时应直接按真正的守恒荷组织巨 ensemble。

若各内部组分分别守恒、等质量且等能，把所有逸度最后统一设为 \(z_a=z\)，则
\[
 \frac1{\mathcal V}\log\Xi(z,\ldots,z)
 =\sum_{n\ge1}
 \left(\sum_{|\mathbf n|=n}\mathfrak c_{\mathbf n}\right)z^n.
\]
与单逸度约定 \eqnrefs{eq:phys-sm-cluster-convention} 比较得到
\begin{equation}
 \frac{b_n}{\lambda_T^3}
 =\sum_{|\mathbf n|=n}\mathfrak c_{\mathbf n}.
 \label{eq:phys-sm-multicomponent-to-single}
\end{equation}
特别在二阶，\(a=b\) 的系数来自 \(\mathfrak c_{2\mathbf e_a}\)，而 \(a\neq b\) 来自 \(\mathfrak c_{\mathbf e_a+\mathbf e_b}\)；后者没有额外的 \(1/2!\)，因为多变量 Taylor 单项式本身已经区分了两个守恒组分。

考虑一个 \(a\)-粒子和一个 \(b\)-粒子的二体扇区。定义总质量和约化质量
\begin{equation}
 m_{ab}^{\rm tot}=m_a+m_b,
 \qquad
 m_{r,ab}=\frac{m_am_b}{m_a+m_b}.
 \label{eq:phys-sm-two-body-masses}
\end{equation}
平移不变二体哈密顿量分离为
\begin{equation}
 H_{ab}
 =\frac{\mathbf P^2}{2m_{ab}^{\rm tot}}
 +H_{{\rm rel},ab},
 \qquad
 H_{{\rm rel},ab}
 =\frac{\mathbf p^2}{2m_{r,ab}}+V_{ab}.
 \label{eq:phys-sm-two-body-separation}
\end{equation}
质心热迹为
\begin{equation}
 Q_{{\rm CM},ab}
 =\frac{\mathcal V}
 {\lambda_T(m_{ab}^{\rm tot})^3}.
 \label{eq:phys-sm-qcm}
\end{equation}
记二体单项式对应的系数为
\[
 \mathfrak c_{ab}^{(2)}:=
 \begin{cases}
 \mathfrak c_{2\mathbf e_a},&a=b,\\
 \mathfrak c_{\mathbf e_a+\mathbf e_b},&a\neq b.
 \end{cases}
\]
因此相互作用引起的二体连通系数改变量满足
\begin{equation}
 \Delta\mathfrak c_{ab}^{(2)}
 =\frac{1}{\lambda_T(m_{ab}^{\rm tot})^3}
 \Delta\mathcal T_{{\rm rel},ab},
 \label{eq:phys-sm-multicomponent-bu-prefactor}
\end{equation}
其中 \(\Delta\mathcal T_{{\rm rel},ab}\) 是相对运动 interacting/自由热迹差，并已经在正确的内部与交换子空间中取迹。

对单组分等质量气体，\(m_{aa}^{\rm tot}=2m\)，而
\[
 \lambda_T(2m)=\frac{\lambda_T}{\sqrt2}.
\]
又有 \(\mathfrak c_2=b_2/\lambda_T^3\)，所以\eqnrefs{eq:phys-sm-multicomponent-bu-prefactor} 特殊化为
\begin{equation}
 \Delta b_2
 =2^{3/2}\Delta\mathcal T_{\rm rel}.
 \label{eq:phys-sm-delta-b2-delta-q2}
\end{equation}
所有常见的 \(\sqrt2\) 因子都来自这个质心高斯归一化，而不是额外的统计简并度。
